%___________________________________________________________________

\begin{frame}{Steuerungstechnik (Nenninger)}
\begin{columns}[T, totalwidth=\textwidth]

\begin{column}{0.5\textwidth}
    
\textbf{Inhalte}
\begin{itemize}
    \item SPS Programmierung
    \item Zustandsgraphen und Schaltnetze
    \item Kodiersysteme
    \item Integration von Sensoren und Aktoren
\end{itemize}

\vspace{0.3cm}
\textbf{Rahmendaten}
\begin{description}
    \item[Klausur] leicht
    \item[Vorlesung] leicht 
    \item[Labor] zeitaufwendig
\end{description}

\end{column}

\begin{column}{0.45\textwidth}
\begin{figure}
    \centering
     \includegraphics[width=\imagewidth, height=0.8\imageheight, keepaspectratio]{siemens}
    \caption{S7 Anlage Dosensortierer.}
\end{figure}
\end{column}

\end{columns}

\end{frame}


%___________________________________________________________________

\begin{frame}{Automatisierungstechnik (Nenninger)}
\begin{columns}[T, totalwidth=\textwidth]

\begin{column}{0.5\textwidth}

\textbf{Inhalte}
\begin{itemize}
    \item Modellbildung technischer Prozesse
    \item Grafische, mathematische Modelle
    \item Feldbus- und Kommunikationssysteme
    \item Zuverlässigkeit, Sicherheit und Verfügbarkeit
    \item Prozessführung, SCADA
\end{itemize}

\vspace{0.3cm}
\textbf{Rahmendaten}
\begin{description}
    \item[Klausur] leicht
    \item[Vorlesung] leicht
    \item[Labor] enorm aufwendig
\end{description}

\end{column}

\begin{column}{0.45\textwidth}
\begin{figure}
    \centering
     \includegraphics[width=\imagewidth, height=0.95\imageheight, keepaspectratio]{ballon}
    \caption{Station Luftballon.}
\end{figure}
\end{column}

\end{columns}

\end{frame}


%___________________________________________________________________

\begin{frame}{Digitale Regelungssysteme (Feßler)}

\textbf{Inhalte}
\begin{itemize}
    \item Digitale Regelungssysteme im zeitdiskreten und $z$-Bereich
    \item Entwurf und Erweiterung von PID-Reglern
    \item Robustheit, Sensitivität, Wasserbett-Effekt
    \item Modellgestützte Verfahren: IMC, MPC, Youla-Parametrierung
    \item Mehrschleifige Regelungen und Prozessführung
    \item Fuzzy Control und digitale Signalverarbeitung
    \item Implementierung auf Signalprozessoren
\end{itemize}

\vspace{0.3cm}
\textbf{Rahmendaten}
\begin{description}
    \item[Klausur] sehr schwer
    \item[Vorlesung] schwer
    \item[Labor] ?
\end{description}

\end{frame}


%___________________________________________________________________

\begin{frame}{Digitale Signalverarbeitung (Strohrmann)}
\begin{columns}[T, totalwidth=\textwidth]

\begin{column}{0.5\textwidth}

\textbf{Inhalte}
\begin{itemize}
    \item Abtasttheorem und diskrete Signale
    \item z-Transformation und Fourier-Analyse
    \item FIR- und IIR-Filterentwurf
    \item DFT und FFT-Algorithmen
    \item Echtzeit-Signalverarbeitung auf DSPs
\end{itemize}

\vspace{0.3cm}
\textbf{Rahmendaten}
\begin{description}
    \item[Klausur] mittel
    \item[Vorlesung] mittel 
    \item[Labor] mittel
\end{description}

\end{column}

\begin{column}{0.45\textwidth}
\begin{figure}
    \centering
     \includegraphics[width=\imagewidth, height=0.8\imageheight, keepaspectratio]{signal}
    \caption{Signaltransformationen.}
\end{figure}
\end{column}

\end{columns}

\end{frame}


%___________________________________________________________________

\begin{frame}{Robotik (Braun)}
\begin{columns}[T, totalwidth=\textwidth]

\begin{column}{0.5\textwidth}

\textbf{Inhalte}
\begin{itemize}
    \item Kinematik und Koordinatentransformationen
    \item Bahnplanung und Bewegungsprofile
    \item Sensorik und Roboterarchitekturen
    \item Programmierung von Robotersystemen
\end{itemize}

\vspace{0.3cm}
\textbf{Rahmendaten}
\begin{description}
    \item[Klausur] leicht
    \item[Vorlesung] leicht 
    \item[Labor] leicht
\end{description}

\end{column}

\begin{column}{0.45\textwidth}
\begin{figure}
    \centering
     \includegraphics[width=\imagewidth, height=0.95\imageheight, keepaspectratio]{roboter}
    \caption{Ein Roboter.}
\end{figure}
\end{column}

\end{columns}

\end{frame}


%___________________________________________________________________

\begin{frame}{Machine Learning in der Automatisierungstechnik (Nenninger)}
\begin{columns}[T, totalwidth=\textwidth]

\begin{column}{0.5\textwidth}

\textbf{Inhalte}
\begin{itemize}
    \item Prozessleitsysteme
    \item Prozesskomponenten SCADA, MES
    \item Datenspeicherung und Cloud-/Datenbankanbindung
    \item Schnittstellen: MQTT, OPC, OPC UA
    \item Grundlagen des Machine Learning für Prozessdaten
\end{itemize}

\vspace{0.3cm}
\textbf{Rahmendaten}
\begin{description}
    \item[Klausur] leicht
    \item[Vorlesung] leicht
    \item[Labor] leicht
\end{description}

\end{column}

\begin{column}{0.45\textwidth}
\begin{figure}
    \centering
     \includegraphics[width=\imagewidth, height=0.95\imageheight, keepaspectratio]{betain}
    \caption{Prozessmodell Betainpumpe.}
\end{figure}
\end{column}

\end{columns}
\end{frame}